% Generated draft: requires mathematical and source review.
\documentclass[11pt]{article}
\usepackage[margin=25mm]{geometry}
\usepackage{fontspec}
\setmainfont{Latin Modern Roman}
\usepackage{amsmath,amssymb,mathtools,unicode-math}
\setmathfont{Latin Modern Math}
\usepackage{enumitem,needspace,xurl,hyperref}
\setlength{\parindent}{0pt}
\setlength{\parskip}{5pt}
\setlength{\emergencystretch}{3em}
\newcommand{\N}{\mathbb{N}}
\newcommand{\Z}{\mathbb{Z}}
\newcommand{\Q}{\mathbb{Q}}
\newcommand{\R}{\mathbb{R}}
\newcommand{\C}{\mathbb{C}}
\newcommand{\F}{\mathbb{F}}
\newcommand{\A}{\mathbb{A}}
\newcommand{\PP}{\mathbb{P}}
\newcommand{\E}{\mathbb{E}}
\newcommand{\eps}{\varepsilon}
\newcommand{\dd}{\,\mathrm{d}}
\begin{document}
\section*{UnsolvedMath ID: 30000478}
\textbf{Source problem code:} catalogue entry
\subsection{Definitions and mathematical statement}
The source's Topological Mirror Symmetry Conjecture says that every family $M$ of Calabi--Yau threefolds has a mirror family $W$ with
\[
 h^{1,1}(W)=h^{2,1}(M),\qquad
 h^{2,1}(W)=h^{1,1}(M),
\]
and with the corresponding Type-II prepotentials identified. For a smooth compact complex fiber $X$, $h^{p,q}(X)=\dim_\C H^q(X,\Omega_X^p)$; family Hodge numbers refer to the fibers in the intended smooth connected family setting.

The usual mirror interpretation exchanges complex-structure deformation data with complexified K\"ahler data and identifies Type-IIA prepotential data on one side with Type-IIB data on the other under an appropriate mirror map. The source requires that correspondence in addition to the numerical Hodge exchange. It does not supply bases of periods, moduli coordinates, prepotential normalizations, the mirror map, or the equivalence under which prepotentials count as identified; prepotentials have coordinate and polynomial ambiguities, so literal equality of unnormalized formulas is not a complete target.

The category of allowed families and mirrors is also unspecified: smooth projective, compact K\"ahler, singular/orbifold, or generalized mirrors are different scopes. This matters for the universal quantifier. If a family has $h^{2,1}(M)=0$ and $W$ is required to be a positive-dimensional smooth compact K\"ahler threefold, the first identity would force $h^{1,1}(W)=0$, whereas a K\"ahler class gives a nonzero $(1,1)$ class. Thus the literal all-family statement requires a scope qualification or a broader mirror notion; it is not silently restricted to a convenient family class.

The literature assessment checked on 21 August 2026 reports Hodge identities and arithmetic relations for important hypersurface and orbifold families while denying a general arithmetic interpretation theorem. Those arithmetic special cases do not by themselves establish this universal complex-geometric and prepotential assertion. The partially-solved label is dated supplied triage requiring independent verification.
\subsection{Short English statement}
Every intended Calabi--Yau threefold family should have a mirror exchanging the two Hodge numbers and matching Type-II prepotential data. The allowed mirror category and normalization are essential, especially for families with no complex deformations.
\subsection{Sources}
\begin{itemize}
\item \emph{Mini-Workshop: Zeta Functions, Index and Twisted K-Theory; Interactions with Physics} (2007), Oberwolfach Reports, DOI 10.4171/owr/2006/21, pp. 1245--1284, \url{https://doi.org/10.4171/owr/2006/21}; proposed year supplied as 2006.
\item Literature assessment checked 21 August 2026: original report, \url{https://doi.org/10.4171/owr/2006/21}; Wan, \emph{Mirror Symmetry for Zeta Functions}, \url{https://arxiv.org/abs/math/0411464}.
\item Local UnsolvedMath registry, permanent ID $30000478$, exact source code \texttt{catalogue entry}, created 25 August 2026 and updated 16 September 2026. Corrected-and-verified status preserves the source-level universal family, Hodge exchange and prepotential claim.
\end{itemize}
\end{document}
